\documentclass[11pt]{article}
\usepackage[a4paper,margin=18mm]{geometry}
\usepackage{fontspec}
\setmainfont{Latin Modern Roman}
\usepackage{amsmath,amssymb,amsthm,amscd,graphicx,color}
\usepackage{xurl}
\usepackage[unicode,hidelinks]{hyperref}
\allowdisplaybreaks
\setlength\emergencystretch{3em}

%\usepackage{amsmath,amsthm,amssymb,amscd}
%\usepackage[dvips]{graphicx}
\usepackage[all]{xy}
%\usepackage[numbers,sort&compress]{natbib}
\usepackage[mathscr]{eucal}
%\usepackage{showkeys}

%%%%%%%%%%%%%%%%%%%special symbols%%%%%%%%%%%%%%%%%%%%%%%%%

\newcommand{\shortleftarrow}{\leftarrow}
\newcommand{\slar}{\shortleftarrow}
%\newcommand{\slar}{\leftarrow}

%%%%%%%%%%%%%%%%%%%%%%itemize%%%%%%%%%%%%%%%%%%%%%%%%%%%%%%

\newenvironment{enm}%
{\begin{enumerate}
\setlength{\itemsep}{0em}
}%
{\end{enumerate}}
\newenvironment{itm}%
{\begin{itemize}
\setlength{\itemsep}{0em}
}%
{\end{itemize}}
\newenvironment{dsc}%
{\begin{description}
\setlength{\itemsep}{0em}
}%
{\end{description}}


\numberwithin{equation}{section}

\newtheorem{Theorem}{Theorem}[section]
\newtheorem{Proposition}[Theorem]{Proposition}
\newtheorem{Lemma}[Theorem]{Lemma}
\theoremstyle{definition}
\newtheorem{Definition}[Theorem]{Definition}
\newtheorem{Remark}[Theorem]{Remark}
\newtheorem{Example}[Theorem]{Example}

\newcommand{\thmref}[1]{Theo\-rem~\ref{#1}}
\newcommand{\secref}[1]{Section~\ref{#1}}
\newcommand{\lemref}[1]{Lemma~\ref{#1}}
\newcommand{\propref}[1]{Proposition~\ref{#1}}
\newcommand{\corref}[1]{Corollary~\ref{#1}}
%\newcommand{\chapref}[1]{Chapter~\ref{#1}}
\newcommand{\dfnref}[1]{Definition~\ref{#1}}
\newcommand{\remref}[1]{Remark~\ref{#1}}
\newcommand{\exref}[1]{Example~\ref{#1}}
\newcommand{\conjref}[1]{Conjecture~\ref{#1}}
\newcommand{\condref}[1]{Condition~\ref{#1}}

%%%%%%%%%%%%%%%%%%%%%%%%MACRO%%%%%%%%%%%%%%%%%%%%%%%%%%%%%

% numbers/spaces
\newcommand{\Z}{\mathbb{Z}} % integers
\newcommand{\Q}{\mathbb{Q}} % rational numbers
\newcommand{\R}{\mathbb{R}} % real numbers
\newcommand{\C}{\mathbb{C}} % complex numbers
\newcommand{\bbP}{\mathbb{P}} % projective space

% Lie groups/algebras
\newcommand{\GL}{\operatorname{GL}} % general linear group
\newcommand{\gl}{\operatorname{\mathfrak{gl}}}
\newcommand{\g}{\mathfrak{g}} % Lie algebra
\newcommand{\h}{\mathfrak{h}} % Cartan subalgebra
\newcommand{\frb}{\mathfrak{b}}
\newcommand{\Lie}{\operatorname{Lie}} % Lie algebra
\newcommand{\Ad}{\operatorname{Ad}} % adjoint
\newcommand{\ad}{\operatorname{ad}}

% linear algebra
\newcommand{\End}{\operatorname{End}}
\newcommand{\Aut}{\operatorname{Aut}}
\newcommand{\Hom}{\operatorname{Hom}}
\newcommand{\Ker}{\operatorname{Ker}}
\newcommand{\range}{\operatorname{Im}}
\newcommand{\Coker}{\operatorname{Coker}}
\newcommand{\codim}{\operatorname{codim}}
\newcommand{\rank}{\operatorname{rank}}
\newcommand{\tr}{\operatorname{tr}}
\newcommand{\unit}{\mathrm{Id}}
\newcommand{\prj}{\operatorname{pr}}

% algebraic geometry
\newcommand{\Spec}{\operatorname{Spec}} % spectrum
%\newcommand{\GIT}{/\hspace{-3pt}/} % good quotient

% quiver
\newcommand{\Rep}{\operatorname{Rep}}
\newcommand{\bM}{\operatorname{\mathbf{M}}}
\newcommand{\bA}{\mathbf{A}} % adjacency matrix
\newcommand{\bD}{\mathbf{D}} % multiplicities
\newcommand{\bC}{\mathbf{C}} % generalized Cartan matrix
\newcommand{\bU}{\mathbf{U}}
\newcommand{\bV}{\mathbf{V}}
\newcommand{\bW}{\mathbf{W}}
\newcommand{\bv}{\mathbf{v}}
\newcommand{\bw}{\mathbf{w}}
\newcommand{\bS}{\mathbf{S}}
\newcommand{\bT}{\mathbf{T}}


% useful tools
\newcommand{\lset}[2]%
{\left\{ \, \left. #1 \hspace{0.25em} \right| \, #2 \, \right\} }
\newcommand{\rset}[2]%
{\left\{ \, #1 \, \left| \hspace{0.25em} #2\right. \, \right\} }
\newcommand{\ov}{\overline}

\newcommand{\res}{\operatorname*{res}}
\newcommand{\mc}{\operatorname{\it mc}}
\newcommand{\add}{\operatorname{\it add}}
\newcommand{\EE}{\operatorname{\mathcal{E}}}
\newcommand{\ved}{\mathbf{d}}
\newcommand{\bZ}{\mathbf{Z}}

\newcommand{\Mirr}{\operatorname{\mathcal{M}}^\mathrm{irr}}
\newcommand{\Mset}{\operatorname{\mathcal{M}}^\mathrm{set}}
\newcommand{\bO}{\mathcal{O}}
\newcommand{\irr}{\mathrm{irr}}
\newcommand{\st}{\mathrm{s}}
\newcommand{\sumfrac}[2]{\genfrac{}{}{0pt}{2}{#1}{#2}}
\newcommand{\hquad}{\hspace{.5em}}


%%%% the first version %%%%%%%%%%%%%
%\newcommand{\Qst}{\operatorname{\mathcal{N}}^\mathrm{s}}
%\newcommand{\Qset}{\operatorname{\mathcal{N}}^\mathrm{set}}
%\newcommand{\vin}{t} % target
%\newcommand{\vout}{s} % source
%\newcommand{\QQ}{\mathcal{Q}}
%\newcommand{\frF}{\mathfrak{F}}

%%%% revised version %%%%%%%%%%%%%%
\newcommand{\Qst}{\operatorname{\mathscr{N}}^\mathrm{s}}
\newcommand{\Qset}{\operatorname{\mathscr{N}}^\mathrm{set}}
\newcommand{\vin}{\mathsf{t}} % target
\newcommand{\vout}{\mathsf{s}} % source
\newcommand{\QQ}{\mathsf{Q}} % quiver
\newcommand{\frF}{\mathscr{F}} % reflection functor

\newcommand{\rd}{\mathop{\mathrm{d}\!}\mathstrut}

\renewcommand{\labelenumi}{{\rm (\theenumi)}}
\renewcommand{\labelenumii}{{\rm (\theenumii)}}
\renewcommand{\theenumi}{\roman{enumi}}

%%%%%%%%%%%%%%%%%%%%%%%%%%%%%%%%%%%%%%%%%%%%%%%%%%%%%%%%%%%%%%%%%%%%%%%%

\begin{document}
\begin{center}\Large For finite loop-free complex quivers with positive vertex multiplicities, reflection at a nonzero leading parameter is an involutive symplectomorphism of stable varieties when both dimension vectors are nonnegative and nonzero\end{center}
\noindent\textbf{Stable ID:} 2821\par
\noindent\textbf{Authors:} Daisuke Yamakawa\par
\noindent\textbf{Original work:} Quiver Varieties with Multiplicities, Weyl Groups of Non-Symmetric Kac-Moody Algebras, and Painlevé Equations\par
\noindent\textbf{Version:} Official SIGMA 6 (2010) article 087, DOI 10.3842/SIGMA.2010.087; journal PDF and native TeX acquired 2026-10-01, exact bytes pinned by SHA256\par
\noindent\textbf{Selected task scope:} Stable-variety part of Theorem4.2 over C. Q finite loop-free, d\_i positive integers; R\_di=\allowbreak{}C[z]/(z\textasciicircum{}di), scalar principal parts lambda\_i(z) in z\textasciicircum{}\{-di\}C[z]/C[z], C=\allowbreak{}2Id-A D. Fix vertex i with lambda\_i,di!=\allowbreak{}0, and nonzero dimension vectors v,s\_i(v) in Z\_\{>=\allowbreak{}0\}\textasciicircum{}I. The original scalar-shift reflection F\_i induces a symplectomorphism N\textasciicircum{}s\_(Q,d)(lambda,v)->N\textasciicircum{}s\_(Q,d)(r\_i(lambda),s\_i(v)) and satisfies F\_i\textasciicircum{}2=\allowbreak{}Id with reflected parameters. r\_i negates lambda\_i and subtracts z\textasciicircum{}\{-1\}c\_ij Res(lambda\_i) from other lambda\_j. No full Weyl-group action on the functors or nonstable quotient claim.\par
\noindent\textbf{Difficulty:} H2: The reflected vertex carries a truncated-polynomial nilpotent module, so the usual kernel/cokernel reflection for ordinary quivers does not provide the required map. The author models the local quotient as a jet-group coadjoint orbit, applies a scalar orbit shift and reconstructs the complementary vertex module. Full moment-map and stability transport calculations, including nilpotent kernel/cokernel arguments, show that this local construction descends to an involutive map preserving the reduced symplectic forms.\par
\noindent\textbf{Editorial boundary note (not author text):} The complete original Theorem4.2 is retained, with selected scope restricted to its stable-variety involutive symplectomorphism and the stated nonzero leading coefficient/nonnegative reflected dimensions. All new truncated-polynomial moment geometry, nilpotent-module stability, jet-group coadjoint-orbit model, scalar shift, complementary vertex reconstruction, moment transport and reduced-form preservation are bundled. The cited prior linear-system orbit theorem remains a prerequisite in the exact author formulation; the current nilpotent kernel/image conditions are included. The original intermediate projection-index notice is preserved. No full Weyl relations, nonstable quotient result, normalization theorem or Painleve assertion is selected.\par
\noindent\textbf{Source:} \url{https://sigma-journal.com/2010/087/}\quad DOI: \url{https://doi.org/10.3842/SIGMA.2010.087}\par
\noindent\textbf{License and attribution:} Copyright retained by the named authors. Published by SIGMA under CC BY 4.0: \url{https://creativecommons.org/licenses/by/4.0/}. Source: \url{https://sigma-journal.com/about.html}. Adaptation consists of standalone excerpt formatting, cover and numbering restoration; original author language and mathematical text are retained.\par
\noindent\textbf{Editorial symbol notice (not author text):} A borrowed general-model line in Lemma3.6 writes V\_hat\_i in an image sum whose target is W; the precise current nilpotent-module kernel/image condition(3.4) is separately given and is used throughout the selected reflection construction. Original formulas unchanged; no reconstruction of the cited earlier theorem. On16 the intermediate moment-change line writes pr\_i(Id\_Vj), while the immediately preceding authored identity is pr\_j(Id\_Vj)=\allowbreak{}d\_j Id\_Vj z\textasciicircum{}\{-1\} and the final result uses c\_ij. Both are retained unchanged; the explicitly authored j-component identity and final parameter formula supply the selected scope, with no collector replacement.\par
\noindent\textbf{Typesetting adaptation (not author text):} The unavailable stmaryrd package is omitted. Its only retained command, shortleftarrow, is rendered with the standard left-arrow glyph for the same incoming-edge labels; all original source calls and arrow direction are unchanged.\par
\medskip\hrule\medskip\noindent\textbf{Author text begins below.}\par
\setcounter{section}{1}
\setcounter{subsection}{0}
\setcounter{subsubsection}{0}
\setcounter{Theorem}{3}
\setcounter{equation}{0}
% Original source lines 520--1698
\section{Preliminaries}\label{sec:pre}

In this section
we brief\/ly recall the def\/inition of Nakajima's quiver variety~\cite{MR1302318}.

\subsection{Quiver}\label{subsec:quiver}

Recall that a (f\/inite) {\em quiver} is a
quadruple $\QQ=(I,\Omega,\vout,\vin)$ consisting of
two f\/inite sets~$I$,~$\Omega$
(the set of {\em vertices}, resp.\ {\em arrows})
and two maps
$\vout, \vin \colon \Omega \to I$
(assigning to each arrow
its {\em source}, resp.\ {\em target}).
Throughout this article, for simplicity,
we assume that {\em our quivers $\QQ$
have no arrow $h \in \Omega$ with $\vout(h)=\vin(h)$}.

For given $\QQ$, we denote by $\ov{\QQ}=(I,\ov{\Omega},\vout,\vin)$
the quiver obtained from $\QQ$ by reversing the orientation of each arrow;
the set $\ov{\Omega}=\{\, \ov h \mid h \in \Omega \,\}$
is just a copy of $\Omega$,
and $\vout(\ov h):=\vin(h)$, $\vin(\ov h):=\vout(h)$
for $h \in \Omega$.
We set $H := \Omega \sqcup \ov{\Omega}$,
and extend the map $\Omega \to \ov{\Omega}$, $h \mapsto \ov h$ to
an involution of $H$ in the obvious way.
The resulting quiver $\QQ + \ov{\QQ} =(I,H,\vout,\vin)$
is called the {\em double} of $\QQ$.

The {\em underlying graph} of $\QQ$,
which is obtained by forgetting the orientation of each arrow,
determines a symmetric matrix $\bA=(a_{ij})_{i,j \in I}$,
called the {\em adjacency matrix}, as follows:
\[
a_{ij} := \sharp\{\,\text{edges joining $i$ and $j$}\,\}
= \sharp\{\, h \in H \mid \vout(h)=i,\ \vin(h)=j \,\}.
\]

Let $\bV=\bigoplus_{i \in I} V_i$ be a nonzero f\/inite-dimensional
$I$-graded $\C$-vector space.
A representation of $\QQ$ over $\bV$ is an element of the vector space
\[
\Rep_\QQ(\bV):= \bigoplus_{h \in \Omega}
\Hom_\C(V_{\vout(h)},V_{\vin(h)}),
\]
and its {\em dimension vector} is given by
$\bv := \dim \bV \equiv (\dim V_i)_{i \in I}$.
Isomorphism classes of representations of $\QQ$ with dimension vector $\bv$
just correspond to orbits in $\Rep_\QQ(\bV)$
with respect to the action of the group
$\GL(\bV):=\prod\limits_{i \in I} \GL_\C(V_i)$ given by
\[
g=(g_i) \colon (B_h)_{h \in \Omega} \longmapsto
\bigl( g_{\vin(h)}B_h g_{\vout(h)}^{-1} \bigr)_{h \in \Omega},
\qquad g \in \GL(\bV).
\]
We denote the Lie algebra of $\GL(\bV)$ by $\gl(\bV)$; explicitly,
$\gl(\bV):=\bigoplus_{i \in I} \gl_\C(V_i)$.
For $\zeta =(\zeta_i)_{i \in I} \in \C^I$,
we denote its image under the natural map $\C^I \to \gl(\bV)$
by $\zeta\,\unit_\bV$,
and also use the same letter $\zeta\,\unit_\bV$
for $\zeta \in \C$ via the diagonal embedding $\C \hookrightarrow \C^I$.
Note that the central subgroup
$\C^\times \simeq
\{\,\zeta\,\unit_\bV \mid \zeta \in \C^\times\,\} \subset \GL(\bV)$
acts trivially on $\Rep_\QQ(\bV)$,
so we have the induced action of the quotient group
$\GL(\bV)/\C^\times$.

Let $B=(B_h)_{h \in \Omega} \in \Rep_\QQ(\bV)$.
An $I$-graded subspace $\bS=\bigoplus_{i \in I}S_i$ of $\bV$ is said to be
{\em $B$-invariant} if
$B_h(S_{\vout(h)}) \subset S_{\vin(h)}$ for all $h \in \Omega$.
If $\bV$ has no $B$-invariant subspace except $\bS=0, \bV$,
then $B$ is said to be {\em irreducible}.
Schur's lemma\footnote{One can apply Schur's lemma
thanks to the following well-known fact:
the category of representations of~$\QQ$
is equivalent to that of an algebra~$\C \QQ$,
the so-called {\em path algebra}; see e.g.~\cite{Etingof}.}
implies that the stabilizer of each irreducible~$B$
is just the central subgroup $\C^\times \subset \GL(\bV)$,
and a standard fact in
Mumford's geometric invariant theory
\cite[Corollary~2.5]{MR1304906} (see also \cite{MR1315461})
implies that the action of $\GL(\bV)/\C^\times$ on the subset
$\Rep_\QQ^\irr(\bV)$ consisting of all irreducible representations
over $\bV$ is proper.

\subsection{Quiver variety}\label{subsec:q-variety}

Suppose that a quiver $\QQ$
and a nonzero f\/inite-dimensional $I$-graded $\C$-vector space
$\bV = \bigoplus_{i \in I}V_i$ are given.
We set
\[
\bM_\QQ(\bV) := \Rep_{\QQ+\ov{\QQ}}(\bV)
= \Rep_\QQ(\bV) \oplus \Rep_{\ov{\QQ}}(\bV),
\]
and regard it as the cotangent bundle of $\Rep_\QQ(\bV)$
by using the trace pairing.
Introducing the function
\[
\epsilon \colon H \to \{ \pm 1 \}, \qquad \epsilon(h):=
\begin{cases}
\hphantom{-}1 & \text{for}\hquad h \in \Omega,\\
-1& \text{for}\hquad h \in \ov{\Omega},
\end{cases}
\]
we can write the canonical symplectic form on $\bM_\QQ(\bV)$ as
\[
\omega := \sum_{h \in \Omega} \tr \rd B_h \wedge \rd B_{\ov{h}}
= \frac12 \sum_{h \in H} \epsilon(h) \tr \rd B_h \wedge \rd B_{\ov{h}},
\qquad (B_h)_{h \in H} \in \bM_\QQ(\bV).
\]
The natural $\GL(\bV)$-action on $\bM_\QQ(\bV)$ is Hamiltonian
with respect to $\omega$ with the moment map
\begin{gather}\label{eq:moment}
\mu=(\mu_i)_{i \in I} \colon \bM_\QQ(\bV) \to \gl(\bV), \qquad
\mu_i(B)=\sum_{\sumfrac{h \in H:}{\vin(h)=i}} \epsilon (h) B_h B_{\ov h}
\end{gather}
vanishing at the origin,
where we identify $\gl(\bV)$
with its dual using the trace pairing.

\begin{Definition}\label{dfn:stable}
A point $B \in \bM_\QQ(\bV)$ is said to be {\em stable}
if it is irreducible as a representation of $\QQ + \ov{\QQ}$.
\end{Definition}

For a $\GL(\bV)$-invariant Zariski closed subset $Z$ of $\bM_\QQ(\bV)$,
let $Z^\st$ be the subset of all stable points in $Z$.
It is a $\GL(\bV)$-invariant Zariski open subset of $Z$,
on which the group $\GL(\bV)/\C^\times$ acts
freely and properly.

\begin{Definition}\label{dfn:q-variety}
For $\zeta \in \C^I$ and $\bv \in \Z_{\geq 0}^I \setminus \{ 0 \}$,
taking an $I$-graded $\C$-vector space $\bV$ with $\dim \bV = \bv$
we def\/ine
\[
\Qst_\QQ(\zeta,\bv) := \mu^{-1}(-\zeta\,\unit_\bV)^\st / \GL(\bV),
\]
which we call the {\em quiver variety}.
\end{Definition}

\begin{Remark}\label{rem:nakajima}
In Nakajima's notation (see \cite{MR1302318} or \cite{MR2023313}),
$\Qst_\QQ(\zeta,\bv)$ is denoted by
$\mathfrak{M}^\mathrm{reg}_{(0,\zeta)}(\bv,0)$.
\end{Remark}



\section{Quiver variety with multiplicities}\label{sec:q-multi}

\subsection{Def\/inition}\label{subsec:dfn}

For a positive integer $d$, we set
\[
R_d := \C[[z]]/z^d\C[[z]], \qquad R^d := z^{-d}\C[[z]]/\C[[z]].
\]
The $\C$-algebra $R_d$ has a typical basis $\{\, z^{d-1}, \dots ,z,1 \,\}$,
with respect to which
the multiplication by $z$ in $R_d$ is represented by
the nilpotent single Jordan block
\[
J_d := \begin{pmatrix}
                        0     & 1 &        & 0 \\
                              & 0 & \ddots &   \\
                              &   & \ddots & 1   \\
                        0     &   &        & 0 \end{pmatrix}
\in \End(\C^d)=\End_\C(R_d).
\]
The vector space $R^d$ may be identif\/ied with
the $\C$-dual $R_d^* =\Hom_\C (R_d,\C)$ of $R_d$ via the pairing
\[
R_d \otimes_\C R^d \to \C, \qquad (f,g) \mapsto
\res_{z=0} \bigl( f(z)g(z) \bigr).
\]

For a f\/inite-dimensional $\C$-vector space $V$, we set
\[
\g_d(V) :=
%\End_{R_d}(V \otimes_\C R_d) \simeq
\gl(V) \otimes_\C R_d = \gl(V)[[z]]/z^d \gl(V)[[z]].
\]
Note that $\g_d(V)$ is naturally isomorphic to
$\End_{R_d}(V \otimes_\C R_d)$ as an $R_d$-module;
hence it is the Lie algebra of the complex algebraic group
\[
G_d(V) := \Aut_{R_d}(V \otimes_\C R_d) \simeq
\lset{g(z) = \sum_{k=0}^{d-1}g_k z^k \in \g_d(V)}{\det g_0 \neq 0}.
\]
The inverse element of $g(z) \in G_d(V)$ is given by
taking modulo $z^d \gl(V)[[z]]$ of
the formal inverse $g(z)^{-1} \in \gl(V)[[z]]$.
The adjoint action of $g(z)$ is described as
\[
(g \cdot \xi)(z)= g(z)\xi(z)g(z)^{-1} \mod z^d \gl(V)[[z]],
\qquad \xi(z) \in \g_d(V).
\]
Using the above $R_d^* \simeq R^d$ and the trace pairing,
we always identify the $\C$-dual $\g_d^*(V)$ of $\g_d(V)$ with
$\gl(V) \otimes_\C R^d = z^{-d}\gl(V)[[z]]/\gl(V)[[z]]$.
Then the coadjoint action of $g(z) \in G_d(V)$ is also described as
\[
(g \cdot \eta)(z)= g(z)\eta(z)g(z)^{-1} \mod \gl(V)[[z]],
\qquad \eta(z)=\sum_{k=1}^d \eta_k z^{-k} \in \g^*_d(V).
\]

The natural inclusion
$\g_d(V) \hookrightarrow \End_\C(V \otimes_\C R_d)
=\End_\C(V) \otimes_\C \End_\C(R_d)$ is represented by
\[
\xi(z)=\sum_{k=0}^{d-1} \xi_k z^k \longmapsto
\sum_{k=0}^{d-1} \xi_k \otimes J_d^k,
\]
whose image is just the centralizer of $\unit_V \otimes J_d$.
Accordingly, its transpose
can be written as
\[
\gl_\C(V \otimes_\C R_d) \simeq \gl_\C (V \otimes_\C R_d)^*
\to \g_d^*(V),\qquad
X \mapsto \sum_{k=1}^d \tr_{R_d}
\bigl[ X \big(\unit_V \otimes J_d^{k-1}\big) \bigr]  z^{-k},
\]
where $\tr_{R_d} \colon \End_\C(V \otimes_\C R_d)
= \End_\C(V)\otimes_\C \End_\C(R_d) \to \End_\C(V)$
denotes the trace of the $\End_\C(R_d)$-part.

Now suppose that a quiver $\QQ$
and a collection of positive integers $\ved =(d_i)_{i \in I} \in \Z_{>0}^I$
are given.
We call the pair $(\QQ,\ved)$ as a {\em quiver with multiplicities}
and $d_i$ as the {\em multiplicity} of the vertex $i$.
Set
\[
R_\ved := \bigoplus_{i \in I} R_{d_i}, \qquad
R^\ved := \bigoplus_{i \in I} R^{d_i},
\]
and for a nonzero f\/inite-dimensional $I$-graded $\C$-vector space
$\bV=\bigoplus_{i \in I} V_i$,
set
\begin{gather*}
\bV_\ved \equiv \bV \otimes_\C R_\ved
:= \bigoplus_{i \in I} V_i \otimes_\C R_{d_i}, \\
\bM_{\QQ,\ved}(\bV) := \bM_\QQ(\bV_\ved)
= \bigoplus_{h\in H}
\Hom_\C \big(V_{\vout(h)} \otimes_\C R_{d_{\vout(h)}},
V_{\vin(h)} \otimes_\C R_{d_{\vin(h)}}\big), \\
G_\ved(\bV) := \prod_{i \in I} G_{d_i}(V_i), \qquad
\g_\ved(\bV):= \bigoplus_{i \in I} \g_{d_i}(V_i).
\end{gather*}
The group $G_\ved(\bV)$ naturally acts on $\bM_{\QQ,\ved}(\bV)$ as a
subgroup of $\GL(\bV_\ved)$.
Note that the subgroup
$\C^\times \subset \GL(\bV_\ved)$
is contained in $G_\ved(\bV)$ and acts trivially on $\bM_{\QQ,\ved}(\bV)$.
As in the case of $\gl(\bV)$,
for $\lambda =(\lambda_i(z))_{i \in I} \in R^\ved$
we denote its image under the natural map
$R^\ved=\g^*_\ved(\C^I) \to \g^*_\ved(\bV)$
by~$\lambda\,\unit_\bV$.

Let $\omega$ be the canonical symplectic form on $\bM_{\QQ,\ved}(\bV)$;
\[
\omega = \frac12 \sum_{h \in H} \epsilon(h) \tr \rd B_h \wedge \rd B_{\ov{h}},
\qquad (B_h)_{h \in H} \in \bM_{\QQ,\ved}(\bV).
\]
Then the $G_\ved(\bV)$-action is Hamiltonian
whose moment map $\mu_\ved$ is given by the composite of
the $\GL(\bV_\ved)$-moment map
$\mu =(\mu_i) \colon \bM_{\QQ,\ved}(\bV) \to \gl(\bV_\ved)$
(see \eqref{eq:moment} for the def\/inition)
and the transpose
$\prj=(\prj_i)$ of the inclusion
$\g_\ved(\bV) \hookrightarrow \gl(\bV_\ved)$;
\begin{gather*}
\mu_\ved  =(\mu_{\ved,i})_{i \in I}
\colon \bM_{\QQ,\ved}(\bV) \to \g_\ved^*(\bV), \\
\mu_{\ved,i}(B)  := \prj_i \circ\, \mu_i (B)
= \sum_{k=1}^{d_i} \sum_{\sumfrac{h \in H:}{\vin(h)=i}}
\epsilon(h) \tr_{R_{d_i}}
\bigl[ B_h B_{\ov{h}} N_i^{k-1} \bigr]  z^{-k},
\end{gather*}
where $N_i := \unit_{V_i} \otimes J_{d_i}$.

\begin{Definition}\label{dfn:conf-stable}
A point $B \in \bM_{\QQ,\ved}(\bV)$
is said to be {\em stable} if
$\bV_\ved$ has no nonzero proper $B$-invariant subspace
$\bS=\bigoplus_{i \in I} S_i$ such that
$S_i \subset V_i \otimes_\C R_{d_i}$
is an $R_{d_i}$-submodule for each $i \in I$.
\end{Definition}

The above stability can be interpreted in terms of
the irreducibility of representations of a~quiver.
Letting $\widetilde{\Omega} := \Omega \sqcup \{ \ell_i \}_{i \in I}$
and extending the maps $\vout$, $\vin$ to $\widetilde{\Omega}$
by $\vout(\ell_i)=\vin(\ell_i)=i$,
we obtain a new quiver
$\widetilde{\QQ}=(I,\widetilde{\Omega},\vout,\vin)$.
Consider the vector space
\[
\Rep_{\widetilde{\QQ}+\ov{\QQ}}(\bV_\ved)
\simeq \bM_{\QQ,\ved}(\bV) \oplus \gl(\bV_\ved)
\]
associated with the quiver
$\widetilde{\QQ}+\ov{\QQ}=(I,\widetilde{\Omega}\sqcup \ov{\Omega},\vout,\vin)$.
Note that in the above def\/inition,
a vector subspace $S_i \subset V_i \otimes R_{d_i}$
is an $R_{d_i}$-submodule
if and only if it is invariant under the action of
$N_i=\unit_{V_i} \otimes J_{d_i}$,
which corresponds to the multiplication by $z$.
Thus letting
\begin{gather}\label{eq:embed}
\iota \colon \ \bM_{\QQ,\ved}(\bV) \hookrightarrow
\Rep_{\widetilde{\QQ}+\ov{\QQ}}(\bV_\ved), \qquad
B \mapsto (B, (N_i)_{i \in I}),
\end{gather}
we see that a point $B \in \bM_{\QQ,\ved}(\bV)$ is
stable if and only if its image $\iota(B)$
is irreducible as a~representation of $\widetilde{\QQ}+\ov{\QQ}$.

For a $G_\ved(\bV)$-invariant Zariski closed subset $Z$ of
$\bM_{\QQ,\ved}(\bV)$,
let $Z^\st$ be the subset of all stable points in $Z$.

\begin{Proposition}
The group $G_\ved(\bV)/\C^\times$ acts
freely and properly on $Z^\st$.
\end{Proposition}

\begin{proof}
Note that
the closed embedding $\iota$ def\/ined in \eqref{eq:embed}
is equivariant under the action of
$G_\ved(\bV) \subset \GL(\bV_\ved)$.
Hence the freeness of
the $G_\ved(\bV)/\C^\times$-action on $Z^\st$
follows from that of the $\GL(\bV_\ved)/\C^\times$-action on
$\Rep_{\widetilde{\QQ}+\ov{\QQ}}(\bV_\ved)^\irr$ and
\[
\iota(\bM_{\QQ,\ved}(\bV)^\st) =
\iota(\bM_{\QQ,\ved}(\bV)) \cap \Rep_{\widetilde{\QQ}+\ov{\QQ}}(\bV_\ved)^\irr,
\]
which we have already checked.
Furthermore, the above implies that the embedding
$Z^\st \hookrightarrow \Rep_{\widetilde{\QQ}+\ov{\QQ}}(\bV_\ved)^\irr$
induced from $\iota$ is closed.
Consider the following commutative diagram:
\[
\xymatrix{
G_\ved(\bV)/\C^\times \times Z^\st
\ar[r] \ar[d] \ar@{}[rd]|{\circlearrowright} & Z^\st \ar[d] \\
\GL(\bV_\ved)/\C^\times \times
\Rep_{\widetilde{\QQ}+\ov{\QQ}}(\bV_\ved)^\irr \ar[r] &
\Rep_{\widetilde{\QQ}+\ov{\QQ}}(\bV_\ved)^\irr,
}
\]
where the vertical arrows are the maps induced from $\iota$,
and the horizonal arrows are the action maps
$(g,x) \mapsto g \cdot x$.
Since the bottle horizontal arrow is proper
and both vertical arrows are closed,
the properness of the top horizontal arrow follows from
well-known basic properties of proper maps
(see e.g.\ \cite[Corollary~4.8]{Hartshorne}).
\end{proof}

\begin{Definition}\label{dfn:q-multi}
For $\lambda \in R^\ved$ and $\bv \in \Z_{\geq 0}^I \setminus \{ 0 \}$,
taking an $I$-graded $\C$-vector space $\bV$ with $\dim \bV = \bv$
we def\/ine
\[
\Qst_{\QQ,\ved}(\lambda,\bv) :=
\mu_\ved^{-1}(-\lambda\,\unit_\bV)^\st / G_\ved(\bV),
\]
which we call the {\em quiver variety with multiplicities}.
\end{Definition}

We also use the following set-theoretical quotient:
\[
\Qset_{\QQ,\ved}(\lambda,\bv) :=
\mu_\ved^{-1}(-\lambda\,\unit_\bV) / G_\ved(\bV).
\]

It is clear from the def\/inition that if $(\QQ,\ved)$ is multiplicity-free,
i.e., $d_i=1$ for all $i \in I$,
then $\Qst_{\QQ,\ved}(\lambda,\bv)$ coincides with the ordinal quiver variety
$\Qst_\QQ(\zeta,\bv)$ with $\zeta_i = \res\limits_{z=0}\lambda_i(z)$.
Even when $\ved$ is non-trivial, for simplicity,
we often refer to $\Qst_{\QQ,\ved}(\lambda,\bv)$ just as the `quiver variety'.

\subsection{Properties}\label{subsec:property}

Here we introduce some basic properties of quiver varieties
with multiplicities.

First, we associate a symmetrizable Kac--Moody algebra
to a quiver with multiplicities $(\QQ,\ved)$.
Let $\bA=(a_{ij})_{i,j \in I}$ be the adjacency matrix of
the underlying graph of $\QQ$ and
set $\bD := (d_i\delta_{ij})_{i,j \in I}$.
Consider the generalized Cartan matrix
\[
\bC=(c_{ij})_{i,j \in I} := 2 \unit - \bA \bD.
\]
Note that it is symmetrizable as $\bD \bC = 2 \bD - \bD \bA \bD$
is symmetric.
Let
\[
\bigl( \g(\bC), \h, \{ \alpha_i \}_{i \in I}, \{ \alpha_i^\vee \}_{i \in I}
\bigr)
\]
be the corresponding Kac--Moody algebra
with its Cartan subalgebra, simple roots and simple coroots.
As usual we set
\[
Q := \sum_{i \in I} \Z \alpha_i, \qquad
Q_+ := \sum_{i \in I} \Z_{\geq 0} \alpha_i.
\]
The diagonal matrix $\bD$
induces a non-degenerate invariant symmetric bilinear form
$(\; , \,)$
on $\h^*$ satisfying
\[
( \alpha_i, \alpha_j) = d_i c_{ij}
= 2 d_i \delta_{ij} - d_i a_{ij} d_j,
\qquad  i,j \in I.
\]

From now on, we regard a dimension vector $\bv \in \Z_{\geq 0}^I$
of the quiver variety as an element of~$Q_+$~by
\[
\Z_{\geq 0}^I \xrightarrow{\simeq} Q_+,
\qquad \bv=(v_i)_{i \in I} \mapsto \sum_{i \in I} v_i \alpha_i.
\]
Let $\res \colon R^\ved \to \C^I$ be the map def\/ined by
\[
\res \colon \ \lambda=(\lambda_i(z)) \longmapsto
\Bigl( \res_{z=0} \lambda_i(z) \Bigr),
\]
and for $(\bv, \zeta) \in Q \times \C^I$,
let $\bv \cdot \zeta := \sum\limits_{i \in I} v_i \zeta_i$
be the scalar product.

\begin{Proposition}\label{prop:dim}\qquad
\begin{enumerate}\itemsep=0pt
\item[{\rm (i)}] The quiver variety $\Qst_{\QQ,\ved}(\lambda,\bv)$
is a holomorphic symplectic manifold of dimension
$2 -(\bv, \bv)$ if it is nonempty.

\item[{\rm (ii)}] If $\bv \cdot \res \lambda \neq 0$,
then $\Qset_{\QQ,\ved}(\lambda,\bv) = \varnothing$.

\item[{\rm (iii)}] If two quivers $\QQ_1$, $\QQ_2$ have the same underlying graph,
then the associated quiver varieties
$\Qst_{\QQ_1,\ved}(\lambda,\bv)$, $\Qst_{\QQ_2,\ved}(\lambda,\bv)$
are symplectomorphic to each other.
\end{enumerate}
\end{Proposition}

\begin{proof}
(i) Assume that $\Qst_{\QQ,\ved}(\lambda,\bv)$ is nonempty.
Since the action of $G_\ved(\bV)/\C^\times$
on the level set $\mu_\ved^{-1}(-\lambda\,\unit_\bV)^\st$ is free and proper,
the Marsden--Weinstein reduction theorem implies that
$\Qst_{\QQ,\ved}(\lambda,\bv)$ is a holomorphic symplectic manifold and
\begin{gather*}
\dim \Qst_{\QQ,\ved}(\lambda,\bv)  =
\dim \bM_{\QQ,\ved}(\bV) - 2 \dim G_\ved(\bV)/\C^\times
 = \sum_{i,j\in I} a_{ij} d_i v_i d_j v_j -2 \sum_{i \in I} d_i v_i^2 + 2 \\
\phantom{\dim \Qst_{\QQ,\ved}(\lambda,\bv)} = {}^t \bv \bD \bA \bD \bv -2 {}^t \bv \bD \bv +2 = 2-(\bv, \bv).
\end{gather*}

(ii) Assume $\Qset_{\QQ,\ved}(\lambda,\bv) \neq \varnothing$
and take a point $[B] \in \Qset_{\QQ,\ved}(\lambda,\bv)$.
Then we have
\[
\sum_{k=1}^{d_i} \sum_{\sumfrac{h \in H:}{\vin(h)=i}}
\epsilon(h) \tr_{R_{d_i}}
\bigl[ B_h B_{\ov{h}} N_i^{k-1} \bigr]\, z^{-k} = -\lambda_i(z)\,\unit_\bV
\]
for any $i \in I$.
Taking $\res\limits_{z=0} \circ \tr$ of both sides and sum over all $i$,
we obtain
\[
\sum_{h \in H} \epsilon(h) \tr ( B_h B_{\ov{h}} )
= -\sum_{i \in I} v_i \res_{z=0} \lambda_i(z) = -\bv \cdot \res \lambda.
\]
Here the left hand side is zero because
\[
\sum_{h \in H} \epsilon(h) \tr ( B_h B_{\ov{h}} ) =
\sum_{h \in H} \epsilon(\ov h) \tr ( B_{\ov{h}}B_h ) =
-\sum_{h \in H} \epsilon(h) \tr ( B_h B_{\ov{h}} ).
\]
Hence $\bv \cdot \res \lambda =0$.

(iii) By the assumption, we can identify the double quivers
$\QQ_1 + \ov{\QQ}_1$ and $\QQ_2 + \ov{\QQ}_2$.
Let $H$ be the set of arrows for them.
Then both the sets of arrows $\Omega_1$, $\Omega_2$ for $\QQ_1$, $\QQ_2$
are subsets of $H$.
Now the linear map
$\bM_{\QQ_1,\ved}(\bV) \to \bM_{\QQ_2,\ved}(\bV)=\bM_{\QQ_1,\ved}(\bV)$
def\/ined by
\[
B \mapsto B', \qquad B'_h:=
\begin{cases}
\hphantom{-}B_h &
\text{if}\hquad h \in \Omega_1 \cap \Omega_2
\hquad\text{or}\hquad h \in \ov{\Omega}_1 \cap \ov{\Omega}_2,
\\
-B_h & \text{otherwise},
\end{cases}
\]
induces a desired symplectomorphism
$\Qst_{\QQ_1,\ved}(\lambda,\bv) \xrightarrow{\simeq}
\Qst_{\QQ_2,\ved}(\lambda,\bv)$.
\end{proof}

Now f\/ix $i \in I$ and set
\[
\widehat{V}_i := \bigoplus_{\vin(h)=i}
V_{\vout(h)} \otimes_\C R_{d_{\vout(h)}}.
\]
Then using it we can decompose the vector space $\bM_{\QQ,\ved}(\bV)$ as
\begin{gather}\label{eq:decomposition}
\bM_{\QQ,\ved}(\bV) =\Hom(\widehat{V}_i,V_i \otimes_\C R_{d_i})
\oplus \Hom(V_i \otimes_\C R_{d_i}, \widehat{V}_i)
\oplus \bM^{(i)}_{\QQ,\ved}(\bV),
\end{gather}
where
\[
\bM^{(i)}_{\QQ,\ved}(\bV) := \bigoplus_{\vin(h), \vout(h) \neq i}
\Hom(V_{\vout(h)} \otimes_\C R_{d_{\vout(h)}},
V_{\vin(h)} \otimes_\C R_{d_{\vin(h)}}).
\]
According to this decomposition,
for a point $B \in \bM_{\QQ,\ved}(\bV)$ we put
\begin{gather*}
B_{i \slar}  := \bigl( \epsilon(h) B_h \bigr)_{\vin(h)=i}
\in \Hom(\widehat{V}_i,V_i \otimes_\C R_{d_i}), \\
B_{\slar i}  := \bigl( B_{\ov{h}} \bigr)_{\vin(h)=i}
\in \Hom(V_i \otimes_\C R_{d_i},\widehat{V}_i), \\
B_{\neq i}  := \bigl( B_h \bigr)_{\vin(h), \vout(h) \neq i}
\in \bM^{(i)}_{\QQ,\ved}(\bV).
\end{gather*}
We regard these as coordinates for $B$ and
write $B=(B_{i \slar},B_{\slar i},B_{\neq i})$.
Note that the symplectic form can be written as
\begin{gather}\label{eq:decomposeform}
\omega = \tr \rd B_{i \slar} \wedge \rd B_{\slar i}
+ \frac12 \sum_{\vout(h),\vin(h) \neq i}
\epsilon(h) \tr \rd B_h \wedge \rd B_{\ov{h}},
\end{gather}
and also the $i$-th component of the moment map can be written as
\[
\mu_{\ved,i}(B) = \prj_i (B_{i \slar} B_{\slar i})
= \sum_{k=1}^{d_i}
\tr_{R_{d_i}} \bigl[ B_{i \slar}  B_{\slar i} N_i^{k-1} \bigr]  z^{-k}.
\]

\begin{Lemma}\label{lem:stability}
Fix $i \in I$ and
suppose that $B$ satisfies at least one of the following two conditions:
\begin{enm}
\item $B$ is stable and $\bv \neq \alpha_i$;
\item the top coefficient
$\tr_{R_{d_i}}(B_{i \slar}  B_{\slar i} N_i^{d_i-1})$
of $\prj_i(B_{i \slar}B_{\slar i})$ is invertible.
\end{enm}
Then $(B_{i \slar},B_{\slar i})$ satisfies
\begin{gather}\label{eq:stability}
\Ker B_{\slar i} \cap \Ker N_i =0, \qquad
\range B_{i \slar} + \range N_i = V_i \otimes_\C R_{d_i}.
\end{gather}
\end{Lemma}

\begin{proof}
First, assume (i) and set
\begin{gather*}
\bS =\bigoplus_{j\in I} S_j, \qquad
S_j :=
\begin{cases}
\Ker B_{\slar i} \cap \Ker N_i & \text{if}\hquad j=i, \\
0 & \text{if}\hquad j \neq i,
\end{cases} \\
\bT =\bigoplus_{j\in I} T_j, \qquad
T_j :=
\begin{cases}
\range B_{i \slar} + \range N_i & \text{if}\hquad j=i, \\
V_j \otimes_\C R_{d_j} & \text{if}\hquad j \neq i.
\end{cases}
\end{gather*}
Then both $\bS$ and $\bT$ are $B$-invariant and
$N_j(S_j) \subset S_j,\, N_j(T_j) \subset T_j$ for all $j \in I$.
Since $B$ is stable,
we thus have $\bS =0$ or $\bS=\bV_\ved$,
and $\bT=0$ or $\bT=\bV_\ved$.
By the assumption $\bv \neq \alpha_i$ and the def\/initions of $\bS$ and $\bT$,
only the case $(\bS,\bT)=(0,\bV_\ved)$ occurs.
Hence $(B_{i \slar},B_{\slar i})$ satisf\/ies~\eqref{eq:stability}.

Next assume (ii).
Set $A(z)=\sum A_k z^{-k}:=\prj_i(B_{i \slar}B_{\slar i})$ and
\[
\widetilde{A} :=
\begin{pmatrix}
0 & 0 & \cdots & 0 \\
\vdots & \vdots & \ddots & \vdots \\
0 & 0 & \cdots & 0 \\
A_{d_i} & A_{d_i-1} & \cdots & A_1
\end{pmatrix}
\in \End_\C(V_i \otimes_\C \C^{d_i})=\End_\C(V_i \otimes_\C R_{d_i}).
\]
Then we have
\[
\tr_{R_{d_i}}(\widetilde{A}N^{k-1}) = A_k
= \tr_{R_{d_i}}(B_{i \slar}  B_{\slar i}N^{k-1}),
\qquad k=1,2,\dots ,d_i,
\]
i.e., $\widetilde{A}-B_{i \slar}B_{\slar i} \in \Ker \prj_i$.
Here, since the group $G_{d_i}(V_i)$ coincides with the centralizer of $N_i$
in $\GL_\C(V_i \otimes_\C R_{d_i})$, we have $\Ker \prj_i = \range \ad_{N_i}$.
Hence there is $C \in \End_\C(V_i \otimes_\C R_{d_i})$ such that
\[
B_{i \slar}  B_{\slar i} = \widetilde{A} + [N_i,C].
\]
Now note that both $\Ker N_i$ and $\Coker N_i$ are naturally isomorphic to $V_i$,
and the natural injection $\iota \colon \Ker N_i \to V_i \otimes_\C R_{d_i}$ and
projection $\pi \colon V_i \otimes_\C R_{d_i} \to \Coker N_i$
can be respectively identif\/ied with
the following block matrices:
\[
\begin{pmatrix} \unit_V \\ 0 \\ \vdots \\ 0
\end{pmatrix} \colon V_i \to V_i \otimes_\C R_{d_i}, \qquad
\begin{pmatrix}
0 & 0 & \cdots & 0 & \unit_V
\end{pmatrix} \colon V_i \otimes_\C R_{d_i} \to V_i.
\]
Thus we have
\begin{gather}\label{eq:momentrel}
\pi B_{i \slar}  B_{\slar i} \iota =
\pi (\widetilde{A} + [N_i,C]) \iota =
\pi \widetilde{A} \iota = A_{d_i}.
\end{gather}
By the assumption $A_{d_i}$ is invertible.
Hence $\pi B_{i \slar}$ is surjective and $B_{\slar i} \iota$ is injective.
\end{proof}


The following lemma is a consequence of results obtained in
\cite{0911.3863}:
\begin{Lemma}\label{lem:orbit1}
Suppose that the set
\begin{gather*}
Z_i :=
\{\, (B_{i \slar},B_{\slar i}) \in \Hom(\widehat{V}_i,V_i\otimes_\C R_{d_i})
\oplus \Hom(V_i\otimes_\C R_{d_i},\widehat{V}_i) \mid \\
\phantom{Z_i := \{\, }{} \prj_i(B_{i \slar}B_{\slar i})=-\lambda_i(z)\,\unit_{V_i},\quad
\text{\rm $(B_{i \slar}, B_{\slar i})$
satisfies \eqref{eq:stability}}
\,\}
\end{gather*}
is nonempty.
Then the quotient of it modulo the action of $G_{d_i}(V_i)$
is a smooth complex manifold having a
symplectic structure induced from
$\tr \rd B_{i \slar} \wedge \rd B_{\slar i}$,
and is symplectomorphic to
a $G_{d_i}(\widehat{V}_i)$-coadjoint orbit via the map given by
\[
\Phi_i \colon \ (B_{i \slar},B_{\slar i}) \longmapsto
- B_{\slar i}(z-N_i)^{-1}B_{i \slar} \in \g^*_{d_i}(\widehat{V}_i).
\]
\end{Lemma}

\begin{proof}
Take any point $(B_{i \slar},B_{\slar i})$ in the above set
and let $\bO$ be the $G_{d_i}(\widehat{V}_i)$-coadjoint orbit
through $\Phi_i(B_{i \slar},B_{\slar i})$.
By Proposition~4,~(a), Theorem~6 and Lemma~3 in \cite{0911.3863},
there exist
\begin{itm}
\item a f\/inite-dimensional $\C$-vector space $W$;
\item a nilpotent endomorphism $N \in \End(W)$ with $N^{d_i}=0$;
\item a coadjoint orbit $\bO_N \subset (\Lie G_N)^*$
of the centralizer $G_N \subset \GL(W)$ of $N$,
\end{itm}
such that the quotient modulo the natural $G_N$-action of the set
\[
\rset{(Y,X) \in \Hom(\widehat{V}_i,W)
\oplus \Hom(W,\widehat{V}_i)}{%
\begin{aligned}
&\prj_N(YX) \in \bO_N,\\
&\Ker X \cap \Ker N =0, \quad
\range Y + \range N =\widehat{V}_i
\end{aligned}
},
\]
where $\prj_N$ is the transpose of the inclusion
$\Lie G_N \hookrightarrow \gl(W)$,
is a smooth manifold having a symplectic
structure induced from $\tr \rd X \wedge \rd Y$,
and is symplectomorphic to $\bO$ via the map
$(Y,X) \mapsto X(z\,\unit_{W} -N)^{-1}Y$.
Note that if $W=V_i\otimes_\C R_{d_i}$, $N=N_i$ and
$\bO_N$ is a single element $\lambda_i(z)\,\unit_{V_i}$,
then we obtain the result by the coordinate change
$(Y,X)=(B_{i \slar},-B_{\slar i})$.
Indeed, this is the case thanks to Proposition~4,~(c) and
Theorem~6 (the uniqueness assertion) in \cite{0911.3863}.
\end{proof}

Note that in the above lemma, the assumption $Z_i \neq \varnothing$
implies $\dim V_i \leq \dim \widehat{V}_i$;
indeed, if $(B_{i \slar},B_{\slar i}) \in Z_i$,
then $B_{\slar i}|_{\Ker N_i}$ is injective by condition \eqref{eq:stability},
and hence
\[
\dim V_i = \dim \Ker N_i = \rank \left( B_{\slar i}|_{\Ker N_i} \right)
\leq \dim \widehat{V}_i.
\]
The following lemma tells us that if the top coef\/f\/icient
of $\lambda_i(z)$ is nonzero,
then the converse is true and
the corresponding coadjoint orbit can be explicitly described:
\begin{Lemma}\label{lem:orbit2}
Suppose $\dim V_i \leq \dim \widehat{V}_i$ and that the top coefficient
$\lambda_{i,d_i}$ of $\lambda_i(z)$ is nonzero.
Then the set $Z_i$ in Lemma~{\rm \ref{lem:orbit1}} is nonempty
and the coadjoint orbit
contains an element of the form
\[
\Lambda_i(z)= \begin{pmatrix}
\lambda_i(z)\,\unit_{V_i} & 0 \\
0 & 0\,\unit_{V'_i} \end{pmatrix},
\]
where $V_i$ is regarded as a subspace of $\widehat{V}_i$ and
$V'_i \subset \widehat{V}_i$ is a complement of it.
\end{Lemma}

\begin{proof}
% First, suppose that the set $Z_i$ is nonempty
% and take any element $(B_{i \slar},B_{\slar i})$.
% Then condition \eqref{eq:stability} implies that
% $B_{\slar i}|_{\Ker N_i}$ is injective.
% Since $\Ker N_i \simeq V_i$, we obtain
% \[
% \dim V_i = \rank \left( B_{\slar i}|_{\Ker N_i} \right)
% \leq \dim \widehat{V}_i.
% \]
Suppose $\dim V_i \leq \dim \widehat{V}_i$
and that the top coef\/f\/icient of $\lambda_i(z)$ is nonzero. We set
\begin{gather*}
B_{i \slar}  := \begin{pmatrix}
0 & 0\, \\ \vdots & \vdots\, \\ 0 & 0\, \\ \unit_{V_i} & 0\,
\end{pmatrix}
\colon \widehat{V}_i = V_i \oplus V'_i \to V_i \otimes_\C R_{d_i}, \\
B_{\slar i}  := - \begin{pmatrix}
\lambda_{i,d_i}\,\unit_{V_i} & \cdots & \lambda_{i,1}\,\unit_{V_i} \\
0 & \cdots & 0
\end{pmatrix}
\colon V_i \otimes_\C R_{d_i} \to \widehat{V}_i,
\end{gather*}
where $\lambda_{i,k}$ denotes the coef\/f\/icient in $\lambda_i(z)$ of $z^{-k}$.
Then we have
\[
\tr_{R_{d_i}}(B_{i \slar}  B_{\slar i} N_i^{k-1})
= -\lambda_{i,k}\,\unit_{V_i}, \qquad k=1,2, \dots ,d_i,
\]
i.e., $\prj_i(B_{i \slar}  B_{\slar i})=-\lambda_i(z)\,\unit_{V_i}$.
The assumption $\lambda_{i,d_i} \neq 0$ and
\lemref{lem:stability} imply that
$(B_{i \slar},B_{\slar i})$ satisf\/ies \eqref{eq:stability}.
Hence $(B_{i \slar},B_{\slar i}) \in Z_i$.
Moreover we have
\begin{gather*}
\Phi_i(B_{i \slar},B_{\slar i})=
-B_{\slar i}(z-N_i)^{-1}B_{i \slar}  =
-\sum_{k=1}B_{\slar i}N_i^{k-1}B_{i \slar}\,z^{-k} \\
\hphantom{\Phi_i(B_{i \slar},B_{\slar i})=
-B_{\slar i}(z-N_i)^{-1}B_{i \slar}}{}
=\sum_{k=1}
\begin{pmatrix}
\lambda_{i,k}\,\unit_{V_i} & 0 \\
0 & 0\,\unit_{V'_i} \end{pmatrix} z^{-k} = \Lambda_i(z).\tag*{\qed}
\end{gather*}
\renewcommand{\qed}{}
\end{proof}

\section{Ref\/lection functor}\label{sec:reflection}

In this section we construct {\em reflection functors}
for quiver varieties with multiplicities.

\subsection{Main theorem}\label{subsec:main}

Recall that the Weyl group $W(\bC)$ of the Kac--Moody algebra $\g(\bC)$
is the subgroup of $\GL(\h^*)$ generated by the simple ref\/lections
\[
s_i(\beta) := \beta -
\langle \beta, \alpha^\vee_i \rangle \alpha_i
= \beta - \frac{2(\beta,\alpha_i)}{(\alpha_i,\alpha_i)}\alpha_i,
\qquad i \in I,\quad \beta \in \h^*.
\]
The fundamental relations for the generators
$s_i,\, i \in I$ are
\begin{gather}\label{eq:coxeter}
s_i^2 =\unit, \qquad
(s_is_j)^{m_{ij}} =\unit, \qquad i,j \in I,\quad i \neq j,
\end{gather}
where the numbers $m_{ij}$ are determined from $c_{ij}c_{ji}$
as the table below
(we use the convention $r^\infty=\unit$ for any $r$)
$$
{\renewcommand{\arraystretch}{1.2}
\begin{array}{|c|c|c|c|c|c|}
\hline
c_{ij}c_{ji} & 0 & 1 & 2 & 3 & \geq 4 \\ \hline
m_{ij} & 2 & 3 & 4 & 6 & \infty \\ \hline
\end{array}}
$$

We will def\/ine a $W(\bC)$-action on the parameter space
$R^\ved \times Q$ for the quiver variety.
The action on the second component $Q$ is given by just
the restriction of the standard action on $\h^*$, namely,
\[
s_i \colon \ \bv=\sum_{i \in I} v_i \alpha_i \longmapsto
\bv - \langle \bv, \alpha^\vee_i \rangle \alpha_i
=\bv - \sum_{j \in I} c_{ij}v_j \alpha_i.
\]
The action on the f\/irst component $R^\ved$ is unusual.
We def\/ine $r_i \in \GL(R^\ved)$ by
\[
r_i(\lambda) = \lambda' \equiv (\lambda'_j(z)), \qquad
\lambda'_j(z) :=
\begin{cases}
-\lambda_i(z) & \text{if}\hquad j =i, \\
\lambda_j(z) - z^{-1}c_{ij} \res\limits_{z=0} \lambda_i(z) & \text{if}\hquad j \neq i.
\end{cases}
\]

\begin{Lemma}\label{lem:rel}
The above $r_i$, $i \in I$ satisfy relations \eqref{eq:coxeter}.
\end{Lemma}

\begin{proof}
The relations $r_i^2 = \unit$, $i \in I$ are obvious.
To check the relation $(r_ir_j)^{m_{ij}}=\unit$ for $i\neq j$,
f\/irst note that the transpose of $s_i \colon Q \to Q$
relative to the scalar product is given by
\[
{}^t s_i \colon \C^I \to \C^I, \qquad
{}^t s_i (\zeta) = \zeta - \zeta_i\sum_{j \in I} c_{ij} \alpha_j.
\]
Now let $\lambda \in R^\ved$. We decompose it as
\[
\lambda = \lambda^0 + \res (\lambda)\,z^{-1},
\qquad \res \lambda^0 =0.
\]
Then we easily see that
\[
r_i(\res(\lambda)\,z^{-1})={}^t s_i (\res(\lambda))\,z^{-1},
\]
and hence that
\[
(r_ir_j)^{m_{ij}} (\lambda) = (r_ir_j)^{m_{ij}} (\lambda^0)
+ \res (\lambda)\,z^{-1}.
\]
Therefore we may assume that $\res \lambda =0$.
Set $\lambda' \equiv (\lambda'_k(z)) := (r_ir_j)^{m_{ij}}(\lambda)$.
Then we have
\[
\lambda'_k(z) =
\begin{cases}
(-1)^{m_{ij}} \lambda_k(z) & \text{if}\hquad k=i,j, \\
\lambda_k(z) & \text{if}\hquad k \neq i,j.
\end{cases}
\]
If $m_{ij}$ is odd, by the def\/inition we have
$c_{ij}c_{ji}=1$.
In particular, $i \neq j$ and
\[
a_{ij} d_j a_{ji} d_i = c_{ij} c_{ji} =1.
\]
This implies $d_i=d_j=1$ and hence that
$\lambda_i(z) = \lambda_j(z)=0$.
\end{proof}

The main result of this section is as follows:
\begin{Theorem}\label{thm:ref}
Let $\lambda=(\lambda_i(z)) \in R^\ved$ and suppose that
the top coefficient $\lambda_{i,d_i}$ of $\lambda_i(z)$
for fixed $i \in I$ is nonzero.
Then there exists a bijection
\[
\frF_i \colon \ \Qset_{\QQ,\ved}(\lambda,\bv) \xrightarrow{\simeq}
\Qset_{\QQ,\ved}(r_i(\lambda),s_i(\bv))
\]
such that $\frF_i^2 = \unit$
and the restriction gives a symplectomorphism
\[
\frF_i \colon \ \Qst_{\QQ,\ved}(\lambda,\bv) \xrightarrow{\simeq}
\Qst_{\QQ,\ved}(r_i(\lambda),s_i(\bv)).
\]
\end{Theorem}

We call the above map $\frF_i$
the {\em $i$-th reflection functor}.

\subsection{Proof of the main theorem}\label{subsec:proof}

Fix $i \in I$ and
suppose that the top coef\/f\/icient
$\lambda_{i,d_i}$ of $\lambda_i(z)$ is nonzero.
Recall the decomposition~\eqref{eq:decomposition} of $\bM_{\QQ,\ved}(\bV)$:
\[
\bM_{\QQ,\ved}(\bV) =\Hom(\widehat{V}_i,V_i \otimes_\C R_{d_i})
\oplus \Hom(V_i \otimes_\C R_{d_i}, \widehat{V}_i)
\oplus \bM^{(i)}_{\QQ,\ved}(\bV),
\]
and the set $Z_i$ given in \lemref{lem:orbit1}.
\lemref{lem:stability} and the assumption $\lambda_{i,d_i} \neq 0$
imply that any
$B=(B_{i \slar},B_{\slar i},B_{\neq i})
\in \mu_{\ved,i}^{-1}(-\lambda_i(z)\,\unit_{V_i})$
satisf\/ies condition \eqref{eq:stability}.
Thus we have
\[
\mu_{\ved,i}^{-1}(-\lambda_i(z)\,\unit_{V_i})
=Z_i \times \bM^{(i)}_{\QQ,\ved}(\bV).
\]
By \lemref{lem:orbit2},
it is nonempty if and only if
\[
v_i \leq \dim \widehat{V}_i
= \sum_j a_{ij} d_j v_j = 2v_i -\sum_j c_{ij}v_j,
\]
i.e., the $i$-th component of $s_i(\bv)$ is non-negative.
We assume this condition,
because otherwise both
$\Qset_{\QQ,\ved}(\lambda,\bv)$ and
$\Qset_{\QQ,\ved}(r_i(\lambda),s_i(\bv))$ are empty
(since $s_i(\bv) \notin \Z_{\geq 0}^I$).
Fix a $\C$-vector space~$V'_i$ of dimension
$\dim \widehat{V}_i - \dim V_i$ and an identif\/ication
$\widehat{V}_i = V_i \oplus V'_i$.
As the group $G_{d_i}(V_i)$ acts trivially on $\bM^{(i)}_{\QQ,\ved}(\bV)$,
Lemmas~\ref{lem:orbit1} and \ref{lem:orbit2} imply that
\[
\mu_{\ved,i}^{-1}(-\lambda_i(z)\,\unit_{V_i})/G_{d_i}(V_i) =
Z_i/G_{d_i}(V_i) \times \bM^{(i)}_{\QQ,\ved}(\bV)
\simeq \bO \times \bM^{(i)}_{\QQ,\ved}(\bV),
\]
where $\bO$ is the $G_{d_i}(\widehat{V}_i)$-coadjoint orbit
through
\[
\Lambda(z)= \begin{pmatrix}
\lambda_i(z)\,\unit_{V_i} & 0 \\
0 & 0\,\unit_{V'_i} \end{pmatrix}.
\]
Now let us def\/ine an $I$-graded $\C$-vector space $\bV'$
with $\dim \bV'=s_i(\bv)$ by
\[
\bV'=\bigoplus_{j \in I} V'_j, \qquad
V'_j := \begin{cases}
V'_i & \text{if}\hquad j=i, \\
V_j & \text{if}\hquad j\neq i,
\end{cases}
\]
and consider the associated symplectic vector space $\bM_{\QQ,\ved}(\bV')$.
Note that $\widehat{V}'_i =\widehat{V}_i$.
Thus by interchanging the roles of $\bV$ and $\bV'$,
$\lambda_i$ and $-\lambda_i$ in Lemmas~\ref{lem:orbit1} and \ref{lem:orbit2},
we obtain an isomorphism
\[
\mu_{\ved,i}^{-1}(\lambda_i(z)\,\unit_{V'_i})/G_{d_i}(V'_i) \simeq
\bO' \times \bM^{(i)}_{\QQ,\ved}(\bV')=\bO' \times \bM^{(i)}_{\QQ,\ved}(\bV),
\]
where $\bO'$ is the $G_{d_i}(\widehat{V}_i)$-coadjoint orbit through
\[
\begin{pmatrix} 0\,\unit_{V_i} & 0 \\
0 & -\lambda_i(z)\,\unit_{V'_i} \end{pmatrix}
=\Lambda(z)-\lambda_i(z)\,\unit_{\widehat{V}_i},
\]
i.e., $\bO' = \bO - \lambda_i(z)\,\unit_{\widehat{V}_i}$.
Hence the scalar shift
$\bO \xrightarrow{\simeq} \bO - \lambda_i(z)\,\unit_{\widehat{V}_i}$
induces an isomorphism
\[
\widetilde{\frF}_i \colon \
\mu_{\ved,i}^{-1}(-\lambda_i(z)\,\unit_{V_i})/G_{d_i}(V_i) \xrightarrow{\simeq}
\mu_{\ved,i}^{-1}(\lambda_i(z)\,\unit_{V'_i})/G_{d_i}(V'_i),
\]
which is characterized as follows: if
\[
\widetilde{\frF}_i[B]=[B'], \qquad
B=(B_{i \slar},B_{\slar i},B_{\neq i}),\qquad
B' =(B'_{i \slar},B'_{\slar i},B_{\neq i}'),
\]
one has
\begin{gather}
B_{\neq i}  = B_{\neq i}', \label{eq:refrelother}\\
- B'_{\slar i}(z-N'_i)^{-1}B'_{i \slar}
 =- B_{\slar i}(z-N_i)^{-1}B_{i \slar} - \lambda_i(z)\,\unit_{\widehat{V}_i},
\label{eq:shift}
\end{gather}
where $N'_i := \unit_{V'_i} \otimes_\C J_{d_i}
\in \End_\C(V'_i \otimes_\C R_{d_i})$.
Note that
\begin{gather}\label{eq:refstability}
\Ker B'_{\slar i} \cap \Ker N'_i =0, \qquad
\range B'_{i \slar} + \range N'_i = V'_i \otimes_\C R_{d_i}
\end{gather}
by \lemref{lem:stability}.

\begin{Lemma}\label{lem:refmoment}
If $\mu_\ved(B)=-\lambda\,\unit_\bV$,
then $\mu_{\ved}(B')=-r_i(\lambda)\,\unit_{\bV'}$.
\end{Lemma}

\begin{proof}
Let $\lambda'=(\lambda'_j(z)) := r_i(\lambda)$.
The identity $\mu_{\ved,i}(B')=\lambda_i(z)\,\unit_{V'_i}$
is clear from the construction.
We check $\mu_{\ved,j}(B')=-\lambda'_j(z)\,\unit_{V'_j}$ for $j \neq i$.
Taking the residue of both sides of~\eqref{eq:shift}, we have
\[
 B'_{\slar i} B'_{i \slar} =  B_{\slar i} B_{i \slar}
+ \lambda_{i,1}\,\unit_{\widehat{V}_i},
\]
which implies that
\[
\epsilon(h) B'_{\ov{h}} B'_h = \epsilon(h) B_{\ov{h}} B_h
+ \lambda_{i,1}\, \unit_{V_{\vout(h)}} \qquad
\text{if} \quad \vin(h)=i.
\]
On the other hand, \eqref{eq:refrelother} means that
$B'_h = B_h$ whenever $\vin(h), \vout(h) \neq i$.
Thus for $j \neq i$, we obtain
\begin{gather}
\sum_{\vin(h)=j} \epsilon(h) B'_h B'_{\ov{h}}
 = \sum_{h \colon i \to j} \epsilon(h) B'_h B'_{\ov{h}} +
\sum_{\vin(h)=j, \vout(h) \neq i} \epsilon(h) B'_h B'_{\ov{h}} \notag\\
\hphantom{\sum_{\vin(h)=j} \epsilon(h) B'_h B'_{\ov{h}}}{}
= \sum_{h \colon i \to j}
\bigl( \epsilon(h) B_h B_{\ov{h}} - \lambda_{i,1}\,\unit_{V_j} \bigr)
+\sum_{\vin(h)=j, \vout(h) \neq i} \epsilon(h) B_h B_{\ov{h}} \notag\\
\hphantom{\sum_{\vin(h)=j} \epsilon(h) B'_h B'_{\ov{h}}}{}
=\sum_{\vin(h)=j} \epsilon(h) B_h B_{\ov{h}}
- a_{ij}\lambda_{i,1}\,\unit_{V_j}.\label{eq:refmoment2}
\end{gather}
Note that
\[
\prj_j(\unit_{V_j}) = \sum_{k=1}^{d_j} \tr_{R_{d_j}} (N_j^{k-1}) z^{-k}
= d_j \unit_{V_j} z^{-1}.
\]
Therefore the image under $\prj_j$
of both sides of~\eqref{eq:refmoment2} gives
\[
\mu_{\ved,j}(B') = \mu_{\ved,j}(B)
- a_{ij}\lambda_{i,1} \prj_i(\unit_{V_j})
=\mu_{\ved,j}(B)
+ c_{ij}\lambda_{i,1}\,\unit_{V_j}z^{-1}.
\]
The result follows.
\end{proof}

\begin{Lemma}\label{lem:refirred}
If $B$ is stable, then so is~$B'$.
\end{Lemma}

\begin{proof}
Suppose that there exists a $B'$-invariant subspace
$\bS'=\bigoplus_j S'_j {\subset} \bV'_\ved$ such that
$N'_j(S'_j) {\subset} S'_j$.
We def\/ine an $I$-graded subspace $\bS=\bigoplus_j S_j$ of
$\bV_\ved$ by
\[
S_j := \begin{cases}
\displaystyle \sum\limits_{k=1}^{d_i} N_i^{k-1} B_{i \slar}\bigl( \widehat{S}'_i \bigr)
& \text{if}\hquad j=i,\\
S'_j & \text{if}\hquad j \neq i,
\end{cases}
\]
where $\widehat{S}'_i := \bigoplus_{\vin(h)=i} S'_{\vout(h)} = \widehat{S}_i$.
Then $B_{i \slar}\bigl( \widehat{S}_i \bigr) \subset S_i$ and
\begin{gather*}
B_{\slar i} (S_i)  = \sum_{k=1}^{d_i}
B_{\slar i} N_i^{k-1} B_{i \slar}\bigl( \widehat{S}'_i \bigr)
 =\sum_{k=1}^{d_i}
\bigl( B'_{\slar i} (N'_i)^{k-1} B'_{i \slar} -\lambda_{i,k} \bigr)
\bigl( \widehat{S}'_i \bigr)
 \subset \widehat{S}'_i = \widehat{S}_i.
\end{gather*}
Hence $\bS$ is $B$-invariant.
Clearly $N_j(S_j) \subset S_j$ for all $j \in I$.
Therefore the stability condition for~$B$ implies that
$\bS=0$ or $\bS=\bV_\ved$.
First, assume $\bS=0$.
Then $S'_j=S_j=0$ for $j \neq i$, and hence
$B'_{\slar i}(S'_i) \subset \widehat{S}'_i=0$, i.e.,
$S'_i \subset \Ker B'_{\slar i}$.
If $S'_i$ is nonzero, then the kernel of the restriction
$N'_i |_{S'_i}$ is nonzero because it is nilpotent.
However it implies $\Ker B'_{\slar i} \cap \Ker N'_i \neq 0$,
which contradicts to \eqref{eq:refstability}.
Hence $S'_i=0$.
Next assume $\bS=\bV_\ved$.
Then $S'_j=S_j=V_j \otimes_\C R_{d_j}$ for $j \neq i$, and hence
$S'_i \supset B'_{i \slar}\bigl( \widehat{S}'_i \bigr) = \range B'_{i \slar}$.
If $V'_i/S'_i$ is nonzero,
then the endomorphism of $V'_i/S'_i$ induced from $N'_i$
has a nonzero cokernel because it is nilpotent.
However it implies
$\range B'_{i \slar} + \range N'_i \neq V'_i \otimes_\C R_{d_i}$,
which contradicts to \eqref{eq:refstability}.
Hence $S'_i=V'_i \otimes_\C R_{d_i}$.
\end{proof}


\begin{proof}[Proof of \thmref{thm:ref}]
As the map $\widetilde{\frF}_i$ is clearly
$\prod\limits_{j \neq i}G_{d_j}(V_j)$-equivariant,
\lemref{lem:refmoment} implies
that it induces a bijection
\[
\frF_i \colon \  \Qset_{\QQ,\ved}(\lambda,\bv) \to \Qset_{\QQ,\ved}(r_i(\lambda),s_i(\bv)),
\qquad [B] \mapsto [B'],
\]
which preserves the stability by \lemref{lem:refirred}.
We easily obtain the relation $\frF_i^2 = \unit$
by noting that $\frF_i$ is induced from the scalar shift
$\bO \to \bO'=\bO -\lambda_i(z)\,\unit_{\widehat{V}_i}$
and the $i$-th component of~$r_i(\lambda)$ is~$-\lambda_i$.
Consider the restriction
\[
\frF_i \colon \ \Qst_{\QQ,\ved}(\lambda,\bv) \to \Qst_{\QQ,\ved}(r_i(\lambda),s_i(\bv)),
\qquad [B] \mapsto [B'].
\]
By \lemref{lem:orbit1} and \eqref{eq:shift},
we have
\[
\tr \rd B_{i \slar} \wedge \rd B_{\slar i} =
\tr \rd B'_{i \slar} \wedge \rd B'_{\slar i},
\]
because the scalar shift
$\bO \to \bO'$
is a symplectomorphism.
Substituting it and \eqref{eq:refrelother}
into \eqref{eq:decomposeform},
we see that the above map $\frF_i$
is a symplectomorphism.
\end{proof}

\begin{Remark}\label{rem:original-ref}
It is clear from \eqref{eq:refrelother}, \eqref{eq:shift}
and \eqref{eq:refstability} that if $d_j=1$ for all $j \in I$,
then $\frF_i$ coincides with the original $i$-th ref\/lection functor
for quiver varieties
(see conditions (a), (b1) and (c) in \cite[Section 3]{MR1623674}).
\end{Remark}

\begin{Remark}\label{rem:ref-relation}
It is known (see e.g.\ \cite{MR1775358}) that if $d_i=1$ for all $i \in I$,
then the ref\/lection functors~$\frF_i$ satisfy relations~\eqref{eq:coxeter}.
We expect that this fact is true for any~$(\QQ,\ved)$.
\end{Remark}

\begin{thebibliography}{99}
\bibitem[9]{Etingof}
Etingof P., Golberg O., Hensel S., Liu T., Schwendner A., Udovina E., Vaintrob D.,
Introduction to representation theory,
\href{http://arxiv.org/abs/0901.0827}{arXiv:0901.0827}.

\bibitem[11]{Hartshorne}
Hartshorne R.,
Algebraic geometry,
{\it Graduate Texts in Mathematics}, Vol.~52, Springer-Verlag, New York~-- Heidelberg, 1977.

\bibitem[16]{MR1315461}
King A.D.,
Moduli of representations of f\/inite-dimensional algebras,
\href{http://dx.doi.org/10.1093/qmath/45.4.515}{\emph{Quart.\ J.\ Math.\ Oxford Ser.~(2)}} \textbf{45} (1994),  515--530.

\bibitem[18]{MR1623674}
Lusztig G.,
On quiver varieties,
\href{http://dx.doi.org/10.1006/aima.1998.1729}{\emph{Adv. Math.}} \textbf{136} (1998), 141--182.

\bibitem[19]{MR1775358}
Lusztig G.,
Quiver varieties and Weyl group actions,
\emph{Ann.\ Inst.\ Fourier (Grenoble)} \textbf{50} (2000), 461--489.

\bibitem[20]{MR1304906}
Mumford D., Fogarty J., Kirwan F.,
Geometric invariant theory, 3rd ed.,
{\it Ergebnisse der Mathematik und ihrer Grenzgebiete (2)}, Vol.~34, Springer-Verlag, Berlin, 1994.

\bibitem[21]{MR1302318}
Nakajima H.,
Instantons on ALE spaces, quiver varieties, and Kac--Moody algebras,
\href{http://dx.doi.org/10.1215/S0012-7094-94-07613-8}{\emph{Duke Math.~J.}} \textbf{76} (1994), 365--416.

\bibitem[22]{MR2023313}
Nakajima H.,
Ref\/lection functors for quiver varieties and Weyl group actions,
\href{http://dx.doi.org/10.1007/s00208-003-0467-0}{\emph{Math. Ann.}} \textbf{327} (2003), 671--721.

\bibitem[31]{0911.3863}
Yamakawa D.,
Middle convolution and Harnad duality,
\href{http://dx.doi.org/10.1007/s00208-010-0517-3}{\emph{Math. Ann.}}, to appear,
\href{http://arxiv.org/abs/0911.3863}{arXiv:0911.3863}.

\end{thebibliography}
\end{document}
